\documentclass[a4paper, 11pt]{article}

% ── Encodage & langue ──────────────────────────────────────────────
\usepackage[utf8]{inputenc}
\usepackage[T1]{fontenc}
\usepackage[french]{babel}

% ── Mise en page ───────────────────────────────────────────────────
\usepackage[top=2cm, bottom=2cm, left=2.2cm, right=2.2cm]{geometry}
\usepackage{parskip}
\setlength{\parindent}{0pt}

% ── Mathématiques ──────────────────────────────────────────────────
\usepackage{amsmath, amssymb}

% ── Tableaux & couleurs ────────────────────────────────────────────
\usepackage{booktabs}
\usepackage{array}
\usepackage[table]{xcolor}
\definecolor{bleu}{RGB}{30, 90, 180}
\definecolor{vertf}{RGB}{30, 140, 80}
\definecolor{rougef}{RGB}{190, 40, 40}
\definecolor{grisclair}{RGB}{240, 243, 248}
\definecolor{gristitre}{RGB}{55, 65, 85}

% ── Boîtes & encadrés ─────────────────────────────────────────────
\usepackage{tcolorbox}
\tcbuselibrary{skins, breakable, theorems}

\tcbset{
  fichebase/.style={
    enhanced, breakable,
    colback=grisclair,
    colframe=bleu,
    fonttitle=\bfseries\color{white},
    coltitle=white,
    attach boxed title to top left={yshift=-2mm, xshift=6mm},
    boxed title style={colback=bleu, colframe=bleu, rounded corners},
    arc=4pt,
    left=6pt, right=6pt, top=10pt, bottom=6pt,
  },
  alertebox/.style={
    enhanced,
    colback=rougef!8,
    colframe=rougef,
    fonttitle=\bfseries\color{rougef},
    coltitle=rougef,
    boxrule=0.8pt,
    arc=3pt,
  },
  verdebox/.style={
    enhanced,
    colback=vertf!8,
    colframe=vertf,
    fonttitle=\bfseries\color{vertf},
    coltitle=vertf,
    boxrule=0.8pt,
    arc=3pt,
  },
}

% ── Tikz (schéma règle du gamma) ───────────────────────────────────
\usepackage{tikz}
\usetikzlibrary{arrows.meta, positioning}

% ── En-tête / pied de page ─────────────────────────────────────────
\usepackage{fancyhdr}
\pagestyle{fancy}
\fancyhf{}
\renewcommand{\headrulewidth}{0.4pt}
\fancyhead[L]{\small\color{gristitre}\textbf{Chimie — Oxydoréduction}}
\fancyhead[R]{\small\color{gristitre}Fiche méthode}
\fancyfoot[C]{\small\thepage}

% ── Titre ──────────────────────────────────────────────────────────
\usepackage{titling}
\pretitle{\begin{center}\Large\bfseries\color{bleu}}
\posttitle{\end{center}}
\preauthor{}\author{}\postauthor{}
\predate{}\date{}\postdate{}

% ══════════════════════════════════════════════════════════════════
\begin{document}
% ══════════════════════════════════════════════════════════════════

\title{Prédire la spontanéité d'une réaction d'oxydoréduction\\[4pt]
\large Méthode complète — Règle du $\gamma$, Nernst et potentiels réels}
\maketitle
\thispagestyle{fancy}
\vspace{-6pt}
\hrule height 1pt
\vspace{10pt}

% ─────────────────────────────────────────────────────────────────
\begin{tcolorbox}[fichebase, title={Rappels fondamentaux}]

\textbf{Couple oxydant/réducteur :} un oxydant Ox et son réducteur conjugué Red sont liés par la demi-équation :
\[
  \text{Ox} + n\,e^- \;\longrightarrow\; \text{Red}
\]

\textbf{Potentiel standard $E°$} (en V, à 25°C, concentrations = 1 mol/L, pH = 0 par convention électrochimique).

\textbf{Potentiel réel $E$} : calculé par la loi de Nernst en tenant compte des conditions effectives (concentrations, pH).

\end{tcolorbox}

\vspace{8pt}

% ─────────────────────────────────────────────────────────────────
\section*{\color{bleu}Méthode en 5 étapes}

% ── Étape 1 ──────────────────────────────────────────────────────
\begin{tcolorbox}[verdebox, title={Étape 1 — Identifier les couples redox en présence}]
Repère dans l'énoncé toutes les espèces susceptibles de jouer le rôle d'oxydant ou de réducteur.\\
Associe chacune à son couple \textbf{Ox/Red} et note son $E°$ (table des potentiels standard).

\medskip
\textit{Exemple :} Fe solide + MgSO\textsubscript{4}(aq) $\;\longrightarrow\;$ couples à considérer :
\[
  \text{Fe}^{2+}/\text{Fe}, \qquad \text{Mg}^{2+}/\text{Mg}, \qquad \text{SO}_4^{2-}/\text{H}_2\text{SO}_3
\]
\end{tcolorbox}

\vspace{4pt}

% ── Étape 2 ──────────────────────────────────────────────────────
\begin{tcolorbox}[verdebox, title={Étape 2 — Écrire les demi-équations redox équilibrées}]
Équilibre les atomes (avec H\textsubscript{2}O) \textbf{puis} les charges (avec H\textsuperscript{+} en milieu acide, OH\textsuperscript{-} en milieu basique).

\medskip
\begin{align*}
  &\text{Fe}^{2+} + 2e^- \longrightarrow \text{Fe} &\quad& (\text{pas de H}^+)\\[4pt]
  &\text{SO}_4^{2-} + 4\,\text{H}^+ + 2e^- \longrightarrow \text{H}_2\text{SO}_3 + \text{H}_2\text{O} &\quad& (\mathbf{4\,\text{H}^+} \text{ impliqués !})
\end{align*}

\begin{tcolorbox}[alertebox, title={Point clé}]
  \textbf{Si H\textsuperscript{+} (ou OH\textsuperscript{-}) figure dans la demi-équation, le potentiel dépend fortement du pH.}\\
  $\Rightarrow$ L'application de Nernst est \textbf{obligatoire}.
\end{tcolorbox}
\end{tcolorbox}

\vspace{4pt}

% ── Étape 3 ──────────────────────────────────────────────────────
\begin{tcolorbox}[verdebox, title={Étape 3 — Calculer les potentiels réels (loi de Nernst)}]

\textbf{Formule générale à 25°C :}
\[
  \boxed{E = E° - \frac{0{,}06}{n}\,\log\frac{[\text{Red}]}{[\text{Ox}]}}
\]

\textbf{Règle pratique :}
\begin{itemize}
  \item Couple \textbf{sans H\textsuperscript{+}} dans la demi-équation : $E \approx E°$ (Nernst change peu si les concentrations sont proches de 1 mol/L).
  \item Couple \textbf{avec $p$ fois H\textsuperscript{+}} : chaque décade de pH fait varier $E$ de $\dfrac{0{,}06\,p}{n}$ V.
  \item \textbf{Si le pH n'est pas précisé dans l'énoncé $\Rightarrow$ supposer pH = 7 (milieu neutre).}
\end{itemize}

\medskip
\textit{Application} — SO\textsubscript{4}\textsuperscript{2-}/H\textsubscript{2}SO\textsubscript{3} à pH 7 (concentrations 1 mol/L) :
\[
  E = 0{,}17 - \frac{0{,}06}{2}\,\log\frac{[\text{H}_2\text{SO}_3]}{[\text{SO}_4^{2-}][\text{H}^+]^4}
  = 0{,}17 - 0{,}03\times\log\!\left(\frac{1}{(10^{-7})^4}\right)
  = 0{,}17 - 0{,}84 \approx -0{,}67\;\text{V}
\]
Le potentiel \textbf{passe de $+0{,}17$ V (en milieu acide) à $-0{,}67$ V (en milieu neutre)}.

\end{tcolorbox}

\vspace{4pt}

% ── Étape 4 ──────────────────────────────────────────────────────
\begin{tcolorbox}[verdebox, title={Étape 4 — Appliquer la règle du $\gamma$ (gamma) avec les potentiels réels}]

\medskip
\textbf{Classe les couples par potentiel réel croissant (de bas en haut) :}

\begin{center}
\begin{tikzpicture}[>=Stealth, font=\small]
  % Axe vertical
  \draw[thick, ->] (0,-0.3) -- (0,3.2) node[right]{$E$ (V)};

  % Couple bas
  \draw[bleu, very thick] (-0.15,0.6) -- (4.5,0.6);
  \node[left] at (-0.15,0.6) {$E_{\text{bas}}$};
  \node[right, bleu] at (4.6,0.6) {\textbf{Ox\textsubscript{2} / Red\textsubscript{2}} \quad réducteur = \textbf{Red\textsubscript{2}}};

  % Couple haut
  \draw[rougef, very thick] (-0.15,2.2) -- (4.5,2.2);
  \node[left] at (-0.15,2.2) {$E_{\text{haut}}$};
  \node[right, rougef] at (4.6,2.2) {\textbf{Ox\textsubscript{1} / Red\textsubscript{1}} \quad oxydant = \textbf{Ox\textsubscript{1}}};

  % Flèches gamma
  \draw[->, thick, rougef, dashed] (2.2,2.2) -- (2.2,0.6);
  \node[right] at (2.25,1.4) {\color{rougef} sens spontané};
\end{tikzpicture}
\end{center}

\[
  \textbf{Réaction spontanée :}\quad \text{Ox}_1 + \text{Red}_2 \;\longrightarrow\; \text{Red}_1 + \text{Ox}_2
\]

L'oxydant du couple \textbf{supérieur} réagit avec le réducteur du couple \textbf{inférieur}.

\end{tcolorbox}

\vspace{4pt}

% ── Étape 5 ──────────────────────────────────────────────────────
\begin{tcolorbox}[verdebox, title={Étape 5 — Vérifier avec $\Delta E$}]

\[
  \boxed{\Delta E = E_{\text{oxydant}} - E_{\text{réducteur}} > 0 \;\Longrightarrow\; \text{réaction \textbf{spontanée}}}
\]
\[
  \Delta E < 0 \;\Longrightarrow\; \text{réaction \textbf{non spontanée} dans ce sens}
\]

\end{tcolorbox}

\vspace{10pt}

% ─────────────────────────────────────────────────────────────────
\section*{\color{bleu}Organigramme récapitulatif}

\begin{center}
\begin{tikzpicture}[
  node distance=0.55cm and 0.3cm,
  box/.style={draw, rounded corners=3pt, fill=grisclair,
              minimum width=6.5cm, minimum height=0.75cm,
              align=center, font=\small},
  dec/.style={draw, diamond, fill=bleu!12, aspect=3.0, align=center,
              font=\small\bfseries, inner sep=2pt},
  arr/.style={->, thick},
  >=Stealth
]
  \node[box, fill=bleu!20] (s1) {\textbf{1.} Identifier les couples Ox/Red};
  \node[box, below=of s1] (s2) {\textbf{2.} Écrire les demi-équations équilibrées};
  \node[dec, below=of s2] (q1) {H\textsuperscript{+} présent ?};
  \node[box, below left=0.6cm and 1.6cm of q1, fill=vertf!12] (oui)
        {\textbf{Oui} $\Rightarrow$ Appliquer Nernst\\(pH = 7 si non précisé)};
  \node[box, below right=0.6cm and 1.6cm of q1, fill=gray!12] (non)
        {\textbf{Non} $\Rightarrow$ Utiliser $E°$ directement};
  \node[box, below=2.2cm of q1] (s4)
        {\textbf{4.} Classer les potentiels réels — Règle du $\gamma$};
  \node[dec, below=of s4] (q2) {$\Delta E > 0$ ?};
  \node[box, below left=0.6cm and 1.4cm of q2, fill=vertf!20] (spon)
        {\textbf{Réaction SPONTANÉE} \checkmark};
  \node[box, below right=0.6cm and 1.4cm of q2, fill=rougef!15] (nspon)
        {\textbf{Non spontanée} \texttimes};

  \draw[arr] (s1) -- (s2);
  \draw[arr] (s2) -- (q1);
  \draw[arr] (q1) -- node[above, font=\scriptsize]{Oui} (oui);
  \draw[arr] (q1) -- node[above, font=\scriptsize]{Non} (non);
  \draw[arr] (oui) -- (s4);
  \draw[arr] (non) -- (s4);
  \draw[arr] (s4) -- (q2);
  \draw[arr] (q2) -- node[above, font=\scriptsize]{Oui} (spon);
  \draw[arr] (q2) -- node[above, font=\scriptsize]{Non} (nspon);
\end{tikzpicture}
\end{center}

\vspace{8pt}

% ─────────────────────────────────────────────────────────────────
\section*{\color{bleu}Exemples illustratifs}

\rowcolors{2}{grisclair}{white}
\begin{tabular}{>{\bfseries}p{4.2cm} p{2.5cm} p{2.5cm} p{1.8cm} p{3.5cm}}
\toprule
\rowcolor{bleu!85}\color{white}Système &
  \color{white}$E°_{\text{haut}}$ &
  \color{white}$E°_{\text{bas}}$ &
  \color{white}Spontané ? &
  \color{white}Raison \\
\midrule
Fe + Cu\textsuperscript{2+} &
  $+0{,}34$ V \newline (Cu\textsuperscript{2+}/Cu) &
  $-0{,}44$ V \newline (Fe\textsuperscript{2+}/Fe) &
  \textcolor{vertf}{\textbf{Oui}} &
  $\Delta E = 0{,}78$ V $> 0$ \\
Cu + Zn\textsuperscript{2+} &
  $-0{,}44$ V (Fe) &
  $-0{,}76$ V (Zn) &
  \textcolor{rougef}{\textbf{Non}} &
  $\gamma$ inversé \\
Zn + Cu\textsuperscript{2+} &
  $+0{,}34$ V \newline (Cu\textsuperscript{2+}/Cu) &
  $-0{,}76$ V \newline (Zn\textsuperscript{2+}/Zn) &
  \textcolor{vertf}{\textbf{Oui}} &
  $\Delta E = 1{,}10$ V $> 0$ \\
Fe + Mg\textsuperscript{2+} &
  $-0{,}44$ V \newline (Fe\textsuperscript{2+}/Fe) &
  $-2{,}37$ V \newline (Mg\textsuperscript{2+}/Mg) &
  \textcolor{rougef}{\textbf{Non}} &
  Mg\textsuperscript{2+} ne peut pas oxyder Fe \\
Fe + SO\textsubscript{4}\textsuperscript{2-} (pH 7) &
  $-0{,}67$ V (réel) &
  $-0{,}44$ V \newline (Fe\textsuperscript{2+}/Fe) &
  \textcolor{rougef}{\textbf{Non}} &
  Nernst abaisse $E$ en dessous de Fe \\
Fe + SO\textsubscript{4}\textsuperscript{2-} (pH 0) &
  $+0{,}17$ V (réel) &
  $-0{,}44$ V \newline (Fe\textsuperscript{2+}/Fe) &
  \textcolor{vertf}{\textbf{Oui}} &
  En milieu très acide seulement \\
\bottomrule
\end{tabular}

\vspace{10pt}

% ─────────────────────────────────────────────────────────────────
\begin{tcolorbox}[alertebox, title={À retenir absolument}]
\begin{enumerate}
  \item \textbf{E° suffit} si aucun H\textsuperscript{+}/OH\textsuperscript{-} dans les demi-équations.
  \item \textbf{Nernst est obligatoire} dès qu'un H\textsuperscript{+} (ou OH\textsuperscript{-}) apparaît.
  \item \textbf{pH = 7 par défaut} si l'énoncé ne précise pas le milieu.
  \item La règle du $\gamma$ s'applique toujours sur les \textbf{potentiels réels}, jamais aveuglément sur $E°$.
  \item $\Delta E > 0$ $\Leftrightarrow$ réaction spontanée dans le sens direct.
\end{enumerate}
\end{tcolorbox}

\end{document}
